\section{Finite recovery of the launch law and prediction}
\label{sec:two-field-finite-law}
The finite experiment also determines the unknown launch law in a
probability metric.  The additional observations use exactly the two
commands of the geometric experiment.  The argument operates on an
isolated occupation component: its normalization is the law of the
sum of a uniform point in the recovered obstacle and the reflected
launch displacement.  Finite moments separate these two factors.
Their conditioning is stated explicitly below.

Write $\mu_0$ for the canonically translated launch probability
measure, and let $W_1$ denote transportation distance for Euclidean
cost.  The moment argument applies to arbitrary compactly supported
probability measures.  The quantitative geometric hypotheses enter
only through the preceding geometric reconstruction.  All windows
and constants in this section are fixed by those hypotheses.

\subsection{Finite occupation moments}
The scalar $a$ in the next two lemmas is an error tolerance; the
command displacement remains the fixed vector $te_1$.
Fix one complete expanded component $P=C_0+(-A_0)$ in the protected
patch.  Its occupation contribution is
\begin{equation}\label{eq:two-field-component-occupation}
 v_C(x)=\int\one_{C_0}(x+z)\,\dd\mu_0(z),\qquad
 \int_{\R^2}v_C(x)\dd x=|C_0|.
\end{equation}
Consequently $v_C/|C_0|$ is the density of $X-Z$, where $X$ is uniform
on $C_0$, $Z$ has law $\mu_0$, and the variables are independent.
This density exists even when $\mu_0$ is singular.

Choose a fixed rational $R\ge1$ sufficiently large that $C_0$, $A_0$,
$P$, and all component cutoffs below lie in $(-R/2,R/2)^2$.
For a multi-index $\alpha=(\alpha_1,\alpha_2)$ put
$x_R^\alpha=(x_1/R)^{\alpha_1}(x_2/R)^{\alpha_2}$.
A geometric estimate with sufficiently small support error supplies
an explicit Lipschitz cutoff $H$ which is one on $P$ and zero on
all other expanded components.  Its Lipschitz constant and support
are uniformly bounded.  For example, apply a piecewise linear
cutoff to the $\ell^\infty$ distance from an estimated polygon for
$P$, with transition width a fixed small fraction of the component
gap.  The error reserve ensures these properties for the true
components.  This construction requires only coarse accuracy for
the cutoff, and $H$ has rational values at rational points.

\begin{lemma}[Moments from finitely many fixed-command bits]
\label{lem:two-field-finite-moments}
Let $m\ge2$, $0<a<1/4$, and $0<\delta<1/4$.  Conditional on a
valid cutoff $H$, the moments
\[
 I_\alpha=\int H(x)v(x)x_R^\alpha\dd x
          =\int v_C(x)x_R^\alpha\dd x,
                    \qquad |\alpha|\le4m,
\]
can all be estimated with error at most $a$ and probability at least
$1-\delta$, using at most
\begin{equation}\label{eq:two-field-moment-cost}
 N_{\rm mom}\le C m^2a^{-4}
                    \log\frac{Cm}{a\delta}
\end{equation}
attempted bits of the two fixed commands.  An integer ceiling is
understood.  Nominal targets form a rational spatial mesh of size
comparable to $a/m$, and the prefix sites extend its fixed window
only by the fixed length $Mt$.  The observations used for the
moments are shared across all multi-indices.
\end{lemma}
\begin{proof}
Put $\ell\asymp a/m$, with a sufficiently small dyadic choice,
and take every mesh node in a fixed square containing the support
of $H$ with a fixed collar.  There are at most $C\ell^{-2}$ nodes.
At each node $x$, estimate the $2M$ means needed by
\eqref{eq:two-field-prefix-inverse}, each to accuracy $c a$.
The maximum of the empirical prefix sums then gives
$|\widehat v(x)-v(x)|\le C M a$; absorb the fixed $M$ by the
choice of the small constant $c$.  Hoeffding's inequality and a
union bound give this event at all mesh nodes using at most
\[
 C\ell^{-2}a^{-2}\log\frac{C}{\ell\delta}
\]
attempted bits.  The rule continues to be valid conditional on
previous observations, because all bits at this stage are fresh.

We verify the spatial quadrature uniformly in the unknown law.
For a fixed $z$, the function being integrated is
$H(x)x_R^\alpha\one_{C_0}(x+z)$.  Its weight has absolute value
at most a constant and Lipschitz constant at most $C(m+1)$.
The total area of mesh cells meeting the boundary of the translated
convex body is at most $C\ell$.  This follows, for example, by
covering the boundary by its inner and outer parallel bands.
On the remaining cells the weight variation is at most $Cm\ell$,
and their total relevant area is bounded.  Thus the Riemann sum
error is at most $Cm\ell$, uniformly for every $z$ in the fixed
footprint bound.  Integration against $\mu_0$ preserves this bound.
Boundary values of an indicator affect only the boundary cells,
so the argument applies also to an atomic launch law.  It uses
neither pointwise continuity of $v$ nor a modulus for $\mu_0$.

Replacing the exact node values by $\widehat v(x)$ adds at most
$Ca$ to every weighted sum.  Decreasing the constants in the
mesh and sampling tolerances proves the accuracy claim and
\eqref{eq:two-field-moment-cost}.  The same empirical node values
are used for all $\alpha$; there is no multiplication of the
observation count by the number of moments.
\end{proof}

\subsection{Quantitative separation of the convolution factors}
The following elementary estimate records the stability needed for
a positive, finite reconstruction of a probability measure.  It also
makes explicit where the nonalgebraic bound in the final acquisition
law enters.  Moment comparison through polynomial approximation is a
classical deconvolution method; see, for example,
\cite[Section~2.3]{RigolletWeed}.  The proof below includes the
bivariate estimate and the perturbation of the recovered convolution
factor required here.  Coordinates in the lemma are dimensionless.

\begin{lemma}[Finite moment separation]
\label{lem:two-field-moment-separation}
Let $U,\widetilde U,Z,\widetilde Z$ be random vectors with probability
laws supported in $[-1,1]^2$, with $U,Z$ independent and
$\widetilde U,\widetilde Z$ independent.  Suppose the moments of
$U$ and $\widetilde U$, and those of $U-Z$ and
$\widetilde U-\widetilde Z$, differ by at most $a$ at every
multi-index of total degree at most $4m$.  Then
\begin{equation}\label{eq:two-field-moment-stability}
 W_1(\mathcal L(Z),\mathcal L(\widetilde Z))
       \le \frac{C}{m}+a(Cm)^{Cm},\qquad m\ge2,
\end{equation}
where $C$ is numerical.  It is enough to assume the bounds on all
the displayed moments; no density or regularity is required.
\end{lemma}
\begin{proof}
For $|\alpha|=k$ the convolution identity is
\begin{equation}\label{eq:two-field-triangular-moments}
 \E(U-Z)^\alpha
   =\sum_{\beta\le\alpha}\binom{\alpha}{\beta}
        \E U^{\alpha-\beta}(-1)^{|\beta|}\E Z^\beta.
\end{equation}
The coefficient of $\E Z^\alpha$ is $(-1)^k$.  Subtract the
corresponding identity for the tilded variables.  All moments of
$U$ and $\widetilde Z$ have absolute value at most one.  If $D_k$
is the largest discrepancy in a moment of $Z$ of total degree $k$,
then $D_0=0$ and
\[
 D_k\le a(1+2^k)+\sum_{r=0}^{k-1}\binom{k}{r}D_r.
\]
Here the binomial identity
$\sum_{|\beta|=r,\,\beta\le\alpha}\binom{\alpha}{\beta}
=\binom{k}{r}$ accounts for both coordinates.  Induction gives
\begin{equation}\label{eq:two-field-factorial-conditioning}
                         D_k\le 2a\,8^k k!.
\end{equation}
Indeed, the sum of the inductive bounds is at most
$2a\,8^k k!(e^{1/8}-1)$, and the forcing term is bounded by the
remaining part of $2a\,8^k k!$.  Thus the triangular inversion has
an explicit finite conditioning bound.

We include the approximation step.  If $f$ is a one-Lipschitz
function on $[-1,1]^2$, subtract its value at zero, so that
$\|f\|_\infty\le2$.  The function
$g(\theta_1,\theta_2)=f(\cos\theta_1,\cos\theta_2)$ is even
in each variable and is Lipschitz on the two-dimensional torus.
Convolve in each variable with the normalized nonnegative kernel
\[
 J_m(\theta)=c_m
       \left(\frac{\sin(m\theta/2)}{\sin(\theta/2)}\right)^4.
\]
Its trigonometric degree is $2(m-1)$, its integral is one, and
\[
 J_m(\theta)\le C\min\{m,m^{-3}|\theta|^{-4}\},\qquad
                 \int_{-\pi}^{\pi}|\theta|J_m(\theta)\dd\theta
                         \le C/m.
\]
The normalization estimate follows by integrating over
$|\theta|\le c/m$ for the lower bound and the displayed elementary
sine bounds for the upper bound.  These bounds also prove the first
moment estimate by splitting at $1/m$.  The tensor convolution
therefore differs from $g$ uniformly by at most $C/m$ and remains
even in each variable.  It has the form
\[
 P(\cos\theta_1,\cos\theta_2),\qquad
 P(x_1,x_2)=\sum_{i,j\le2m}b_{ij}T_i(x_1)T_j(x_2),
\]
where $T_i$ is the Chebyshev polynomial.  Its cosine coefficients
satisfy $|b_{ij}|\le C$.  The recurrence
$T_{i+1}=2xT_i-T_{i-1}$ bounds the sum of the absolute monomial
coefficients of $T_i$ by $3^i$.  Consequently the corresponding
coefficient sum for $P$ is at most $Cm^2 3^{4m}$.

Applying \eqref{eq:two-field-factorial-conditioning} to this
polynomial, whose total degree is at most $4m$, gives
\[
 |\E f(Z)-\E f(\widetilde Z)|
       \le C/m+a(Cm)^{Cm}.
\]
Taking the supremum over Lipschitz functions proves
\eqref{eq:two-field-moment-stability}, by the transportation duality
on a compact set.  For finite laws this is ordinary finite linear
programming duality; approximation on finer finite grids proves the
same identity for arbitrary compact probability measures.
\end{proof}

\subsection{A finite probability output}
Let $N_{\rm geom}(u,\delta)$ denote the observation bound in
Theorem~\ref{thm:two-field-finite-geometry}, with geometric tolerance
$u$.  In particular it supplies a matched body in the canonical
frame, its expanded component, and the protected patch.  The next
statement separates the two costs before substituting that bound.

\begin{theorem}[Finite recovery of the geometry and launch probability]
\label{thm:two-field-finite-law}
Assume the hypotheses of Theorem~\ref{thm:two-field-finite-geometry}.
At one unknown stationary law, the same two fixed opposite commands
admit a finite controller which returns the geometric output of that
theorem and a finitely supported probability measure $\widehat\mu_0$.
For all sufficiently small $\varepsilon>0$ and $0<\delta<1/4$,
with probability at least $1-\delta$ the geometric loss is at most
$C\varepsilon$, the discrete data are correct, and
\begin{equation}\label{eq:two-field-finite-wasserstein}
                         W_1(\widehat\mu_0,\mu_0)
                                      \le C\varepsilon.
\end{equation}
More precisely, let $m=\lceil C_0/\varepsilon\rceil$ and take a
certified dyadic $a_m$ with
\begin{equation}\label{eq:two-field-moment-tolerance}
 2^{-C_1m\log_2(C_1m)-1}
       \le a_m\le 2^{-C_1m\log_2(C_1m)},
\end{equation}
where $C_0,C_1$ are sufficiently large fixed constants.  The
attempted-bit bound is
\begin{equation}\label{eq:two-field-joint-decomposed-cost}
 N\le N_{\rm geom}(c a_m,\delta/2)
       +Cm^2a_m^{-4}\log\frac{Cm}{a_m\delta}.
\end{equation}
Inserting the polynomial geometric bound gives the explicit uniform
sufficient bound
\begin{equation}\label{eq:two-field-joint-exponential-cost}
 N\le\exp\!\left(C\varepsilon^{-1}
                       \log\frac{C}{\varepsilon}\right)
                    \log^2\frac{C}{\delta},
                \qquad S\le J\le N.
\end{equation}
Here every bit, site visit and fixed-command occurrence is charged.
The two displacements remain $a$ and $-a$; their length does not
shrink with $\varepsilon$.  All nominal sites lie in a fixed
prior-dependent aperture.  The coordinates used by the moment stage
have mesh comparable to $a_m/m$; together with the geometric stage
the command description has at most
\begin{equation}\label{eq:two-field-joint-description}
 CJ\left\{\varepsilon^{-1}\log(C/\varepsilon)
                            +\log(C/\delta)\right\}
\end{equation}
bits.  The same coarse bound on the unknown footprint origin as in
the geometric theorem gives a uniform bound on the absolute physical
aperture.  No evaluation of the density, second preparation setting,
angular derivative, or additional displacement vector is supplied.

The moment stage itself is valid for arbitrary compact launch
probabilities whenever the stated geometric estimates are available.
The exponential bound is a sufficient bound for this two-command
joint inverse; it is not a matched minimax assertion.  Under the
complete finite-configuration hypotheses of
Corollary~\ref{cor:two-field-finite-cloud}, the same acquisition and
probability conclusions hold for all the recovered components,
without a periodicity or nonperiod-patch hypothesis.
\end{theorem}
\begin{proof}
Use half the failure allowance to reconstruct geometry at tolerance
$c a_m$.  Work on its common success event, and take one complete
matched obstacle $\widehat C$ and the rational polygon for its
expanded component.  They give the cutoff $H$ above.  The true body
$C_0$ and $\widehat C$ have Hausdorff distance at most $C c a_m$.
Their areas are uniformly bounded below, and the symmetric
difference of the bodies has area at most $C c a_m$.  The latter
bound follows from the convex parallel-body formula and the uniform
perimeter bound.  Thus the uniform probability measures on the two
bodies have total variation distance at most $C c a_m$.
Every normalized coordinate monomial has absolute value at most one,
so their moments differ by this same bound, independently of degree.
All geometric moments through degree $4m$ can be computed from the
finite geometric representation with certified additional error at
most $c a_m$.  This requires no observations.  For example, uniform
inner and outer rational polygon approximations give area errors and
polynomial integral errors tending to zero with explicit support and
Lipschitz bounds.

Use Lemma~\ref{lem:two-field-finite-moments}, with the remaining
failure allowance, to estimate all $I_\alpha$ to error $c a_m$.
Normalize by a certified positive rational approximation to
$|\widehat C|$, with error at most $c a_m$, to obtain the moments
of $Y=(X-Z)/R$ to error $C c a_m$.  The true moment identities use
$U=X/R$ and $Z/R$, in a fixed dimensionless box; the reconstructed
body supplies the moments of $\widehat U$.  Increase the fixed $R$
if necessary so that all quantities and their reserves are in
$[-1,1]^2$.

We describe an entirely finite positive reconstruction.  Take the
rational grid of spacing $h\asymp c a_m/m$ in $[-1,1]^2$.  It may
be restricted to a known $Ca_m$-parallel neighborhood of the
dimensionless footprint $R^{-1}\widehat A_0$, with its padding larger
than the rescaled certified Hausdorff error plus $2\sqrt2h$. Every
point of $R^{-1}A_0$ then has a nearest retained grid point within $Ch$.  For unknown
nonnegative weights $(p_z)$ summing to one on this finite grid,
form the convolution moments
\[
 T_\alpha(p)=\sum_{\beta\le\alpha}\binom{\alpha}{\beta}
       \widetilde c_{\alpha-\beta}(-1)^{|\beta|}
                         \sum_zp_z z^\beta,
                       \qquad |\alpha|\le4m,
\]
where $\widetilde c$ are the certified moments of $\widehat U$.
Minimize the largest discrepancy between these quantities and the
estimated moments of $Y$.  This is a finite linear program with
rational coefficients; its feasible set is the probability simplex.
A rational solution to any prescribed objective tolerance exists
and can be certified by finite linear programming.

The law of $Z/R$, moved to its nearest retained grid point, is a
comparison probability on this simplex.  Its effect on a monomial
of the difference $\widehat U-Z/R$ of degree at most $4m$ is bounded
by $Cmh$, because all difference coordinates have absolute value
strictly below one after the choice of $R$.  The errors in the
geometric moments contribute at most $2^{4m}c a_m$, if bounded
term by term; this harmless exponential factor can be absorbed
below.  Thus the optimum, with an additional certified tolerance
$a_m$, is at most $C2^{4m}a_m$.
The fitted law and the true law consequently satisfy the assumptions
of Lemma~\ref{lem:two-field-moment-separation} with moment tolerance
at most $C2^{4m}a_m$.  The same factor covers the numerical moment
errors.  The constants in \eqref{eq:two-field-moment-tolerance}
ensure
\[
                  C2^{4m}a_m(Cm)^{Cm}\le 1/m.
\]
After multiplying coordinates by $R$, the fitted law is the claimed
$\widehat\mu_0$, and its transportation error is at most $C/m$.
There are at most $C(m/a_m)^2$ possible output atoms, all with
rational weights and locations.  Numerical construction and the
finite output size are separate from attempted-bit counts.

The union of the two success events has probability at least
$1-\delta$.  Lemma~\ref{lem:two-field-finite-moments} gives
\eqref{eq:two-field-joint-decomposed-cost}.  The polynomial geometric
cost at $c a_m$ and the displayed formula for $a_m$ imply
\eqref{eq:two-field-joint-exponential-cost}; factors polynomial in
$m$ and $\log(1/a_m)$ are absorbed by increasing its fixed constant.
The geometric error $Ca_m$ is smaller than $C\varepsilon$.
Every spatial grid point and its prefix translates use the original
two displacements.  The coordinate denominator has logarithm at
most $C\varepsilon^{-1}\log(C/\varepsilon)$, and the same holds
for the geometric-stage meshes.  Charging repetitions separately
proves the remaining resource and description assertions.  For the
complete finite configuration, the cited corollary provides the same
isolated component and geometric estimates.  Every moment and
probability step above is unchanged and uses no period relation.
\end{proof}

\subsection{Finite prediction of all bounded-length responses}
The probability output has a direct experimental interpretation.
For a bounded nominal window $U$ and a fixed maximum command length
$T$, define the predicted forward family by
\begin{equation}\label{eq:two-field-predicted-family}
 \widehat F_b(x)=\int B_b^{\widehat{\mathcal O}_0}(x+z)
                                      \,\dd\widehat\mu_0(z),
                                  \qquad |b|\le T.
\end{equation}
This is a finite weighted sum for every specified $b$.  The commands
in this predicted family need not have been used in acquisition.

\begin{corollary}[Prediction in local spatial mean]
\label{cor:two-field-finite-prediction}
On the success event of Theorem~\ref{thm:two-field-finite-law},
\begin{equation}\label{eq:two-field-prediction-Lone}
 \sup_{|b|\le T}\int_U
             |\widehat F_b(x)-F_b(x)|\dd x
                         \le C_{U,T}\varepsilon.
\end{equation}
No density upper bound or continuity modulus is required.  More
generally, the right side is $C_{U,T}(\xi+\epsilon)$ if the
matched geometry in the relevant patch has Hausdorff error $\xi$
and $W_1(\widehat\mu_0,\mu_0)\le\epsilon$.  Spatial averages
against any bounded test supported in $U$ inherit this error times
its supremum norm.  The entire predicted family uses the finite
geometric and probability output and needs no further observations.
\end{corollary}
\begin{proof}
Write $D_b(C)=C+[-b,0]$ for the convex swept body.  Up to boundary
sets of planar measure zero, the free-start collision set of a
configuration is
\[
 E_b^{\mathcal O}
   =\left(\bigcup_C D_b(C)\right)\setminus\left(\bigcup_C C\right),
                         \qquad B_b^{\mathcal O}=\one_{E_b^{\mathcal O}}.
\]
Only a uniformly bounded number of bodies meet any fixed relevant
window when $|b|\le T$.  Their swept bodies have perimeters at most
$\operatorname{Per}(C)+2T$.  Hence the collision indicator has a
uniform bound on its total variation on a slightly larger window.
If matched bodies have Hausdorff error at most $\xi$, so do their
swept bodies.  Convex parallel-body area bounds, summed over these
finitely many components, therefore give
\begin{equation}\label{eq:two-field-geometric-indicator-stability}
 \|B_b^{\widehat{\mathcal O}_0}-B_b^{\mathcal O_0}\|_{L^1(V)}
                               \le C_{V,T}\xi.
\end{equation}
For periodic tables, the geometric loss supplies this matching on
every prescribed bounded patch, with a constant depending on that
patch.  No uniform assertion on the whole plane with perturbed
lattice vectors is used.

The translation estimate for a function of bounded variation
(see \cite[Chapter~3]{AFP}) gives
\begin{equation}\label{eq:two-field-BV-translation}
 \int_U |B_b(x+z)-B_b(x+\widetilde z)|\dd x
                              \le C_{U,T}|z-\widetilde z|,
\end{equation}
uniformly for $z,\widetilde z$ in the common bounded footprint box.
For a smooth function this follows by integrating its directional
derivative on segments and applying Fubini.  Approximation in
$L^1$ with convergence of total variations proves the same bound
for the displayed finite-perimeter indicator.  All intermediate
segments stay in a fixed enlarged window, where its variation has
already been bounded.

Integrate the sum of
\eqref{eq:two-field-geometric-indicator-stability} and
\eqref{eq:two-field-BV-translation} against an optimal coupling of
$\mu_0$ and $\widehat\mu_0$, and apply Fubini once more.  This gives
$C_{U,T}(\xi+W_1(\widehat\mu_0,\mu_0))$, uniformly in $b$.
Boundary conventions make no difference to these spatial integrals,
even for an atomic launch law.  The theorem's estimates imply
\eqref{eq:two-field-prediction-Lone}.
\end{proof}

\begin{corollary}[A density estimate under a variation bound]
\label{cor:two-field-finite-density}
Suppose, in addition, that $\mu_0=j_0\dd z$ and that the zero
extension of $j_0$ has total variation at most a known $V$.
There is a finite density output $\widehat j\ge0$, of integral one,
with
\[
                    \|\widehat j-j_0\|_{L^1(\R^2)}
                                    \le C(V+1)\varepsilon
\]
and success probability at least $1-\delta$, using the same two
fixed commands.  A sufficient attempted-bit bound is
\begin{equation}\label{eq:two-field-BV-density-cost}
 \exp\!\left(C\varepsilon^{-2}\log(C/\varepsilon)\right)
                                  \log^2(C/\delta).
\end{equation}
For every fixed bounded nominal window $U$ and $T<\infty$, this
same density output also gives the raw predictions
\begin{align}
 \widehat F_b^{\rm dens}(x)
   &=\int B_b^{\widehat{\mathcal O}_0}(x+z)
                                  \widehat j(z)\dd z,
                       \label{eq:two-field-density-prediction}\\
 \sup_{\substack{x\in U\\ |b|\le T}}
      |\widehat F_b^{\rm dens}(x)-F_b(x)|
   &\le C_{U,T,V}\varepsilon.
                       \label{eq:two-field-uniform-prediction}
\end{align}
The variation hypothesis concerns these strong density and uniform
raw-mean conclusions.  It is not needed for the probability or
local spatial-mean conclusions.
\end{corollary}
\begin{proof}
Apply Theorem~\ref{thm:two-field-finite-law} at transportation
tolerance $c\varepsilon^2$.  Let $\rho$ be a fixed nonnegative,
smooth compactly supported probability kernel, put
$\rho_h(z)=h^{-2}\rho(z/h)$, and output
$\widehat j=\rho_\varepsilon*\widehat\mu_0$.
This is a finite mixture of explicitly specified densities.
The translation estimate for $BV$ functions gives
\[
 \|\rho_h*j_0-j_0\|_1\le CVh.
\]
For any coupling of the two laws,
$\|\rho_h(\cdot-z)-\rho_h(\cdot-\widetilde z)\|_1
\le C h^{-1}|z-\widetilde z|$.
Integrating an optimal coupling therefore gives
\[
 \|\rho_h*\widehat\mu_0-j_0\|_1
          \le C\{Vh+h^{-1}W_1(\widehat\mu_0,\mu_0)\}.
\]
Take $h=\varepsilon$ and substitute the theorem's cost at
$c\varepsilon^2$.

For the uniform raw predictions, first observe the planar $BV$
bound
\[
 \|j_0\|_2
 \le\int_0^\infty |\{j_0>q\}|^{1/2}\dd q
 \le C\int_0^\infty\operatorname{Per}(\{j_0>q\})\dd q
 =C|Dj_0|(\R^2)\le CV.
\]
The first inequality is Minkowski's integral inequality applied to
the layer-cake representation, the second is the planar
isoperimetric inequality, and the equality is the coarea formula.
For fixed $U,T$, all evaluations $x+z$ and their commanded segments
lie in a fixed enlarged canonical window.  Its relevant obstacle
count is bounded by the geometric priors.  The matching statement
and the swept-body estimate
\eqref{eq:two-field-geometric-indicator-stability} therefore imply,
uniformly for $x\in U$ and $|b|\le T$,
\[
 \int j_0(z)
  |B_b^{\widehat{\mathcal O}_0}(x+z)
       -B_b^{\mathcal O_0}(x+z)|\dd z
                 \le C_{U,T}\|j_0\|_2\sqrt\xi
                 \le C_{U,T}V\sqrt\xi,
\]
where $\xi$ is the geometric tolerance.  The bound follows by
Cauchy--Schwarz and the area bound on the indicator's symmetric
difference.  Replacing $j_0$ by $\widehat j$ in the estimated
geometry adds at most $\|\widehat j-j_0\|_1$ to any raw mean.
In the present construction $\xi=Ca_m$ with
$m\asymp\varepsilon^{-2}$, so $\sqrt\xi\le C\varepsilon$.
These bounds prove \eqref{eq:two-field-uniform-prediction}.
The use of a fixed enlarged window also accounts for perturbations
of lattice vectors in a periodic output.
\end{proof}
